\clearpage\section*{附录 E}\phantomsection\addcontentsline{toc}{section}{附录 E}
\label{app:samples}

\subsection*{生成样本}
\vspace{2ex}
\begin{table}[htbp]

\begin{center}
\begin{tabular}[t]{@{}ccc@{}}
 \includegraphics[width=0.25\textwidth]{\figuresdir{samples1vae.jpg}} & \hspace{3ex} &
\includegraphics[width=0.25\textwidth]{\figuresdir{samples1iwae.jpg}} \\
& & \\
 \includegraphics[width=0.25\textwidth]{\figuresdir{samples2vae.jpg}} & \hspace{3ex} &
\includegraphics[width=0.25\textwidth]{\figuresdir{samples2iwae.jpg}} \\
\end{tabular}
\end{center}

\caption{从 VAE（左列）和 $k=50$ 的 IWAE（右列）模型中随机生成的样本。第一行：单随机层模型。第二行：双随机层模型。样本以相应伯努利分布的均值表示。}
\label{tbl:samples}
\end{table}
